%
%
%
%
%
%
%
%
%
%
%
%
%
%
%
%
%
%
%
%
%
%
%
%
%
%
%
%
%
%
%
%
%
%

\usepackage{makeidx}
\makeindex
\renewcommand{\seename}{siehe}
\newif\ifZSFIndexShowPageNumber
\ZSFIndexShowPageNumberfalse

% ------------------------------------------------------------
% Sofort-Writer
% ------------------------------------------------------------
\makeatletter
% Aktuelle Abschnittsnummer einfrieren; wenn vorhanden die dritte Ebene
% (x.x.x), sonst die zweite Ebene (x.x), vor der ersten \SubsectionBar
% eines Kapitels Fallback auf die Kapitelnummer.
\newcommand{\ZSF@idxnumnow}{%
  \ifnum\value{subsubsection}>0 \thesubsubsection\else
    \ifnum\value{subsection}>0 \thesubsection\else\thesection\fi
  \fi}
% #1 = Eintrag (ggf. sort@display), #2 = Encap (ZSFidxnum | see{...}).
% Im Seitenzahl-Modus wird ein zusammengesetzter Locator Abschnitt.Seite
% geschrieben. makeindex kann ihn mit page_compositor "." sortieren und
% verdichten; die Ausgabe trennt beide Teile erst am Schluss wieder.
\newcommand{\ZSF@writeidx}[2]{%
  \@ifundefined{@indexfile}{}{%
    \begingroup
      \edef\ZSF@num{\ZSF@idxnumnow}%
      \ifZSFIndexShowPageNumber
        \edef\ZSF@locator{\ZSF@num.\thepage}%
        \protected@edef\ZSF@line{\string\indexentry{#1|#2}{\ZSF@locator}}%
      \else
        \protected@edef\ZSF@line{\string\indexentry{#1|#2}{\ZSF@num}}%
      \fi
      \immediate\write\@indexfile{\ZSF@line}%
    \endgroup}}
% «siehe»-Verweise bleiben reine Abschnitts-Wegweiser und erhalten nie
% eine Seitenzahl.
\newcommand{\ZSF@writeidxsee}[2]{%
  \@ifundefined{@indexfile}{}{%
    \begingroup
      \edef\ZSF@num{\ZSF@idxnumnow}%
      \protected@edef\ZSF@line{\string\indexentry{#1|see{#2}}{\ZSF@num}}%
      \immediate\write\@indexfile{\ZSF@line}%
    \endgroup}}

% ------------------------------------------------------------
% Öffentliche Eintrags-Makros
% ------------------------------------------------------------
% (bleiben im \makeatletter-Block: Bodies referenzieren \ZSF@writeidx)
\NewDocumentCommand{\ZSFindex}{o m}{%
  \IfNoValueTF{#1}%
    {\ZSF@writeidx{#2}{ZSFidxnum}}%
    {\ZSF@writeidx{#1@#2}{ZSFidxnum}}}
\NewDocumentCommand{\ZSFindexsee}{o m m}{%
  \IfNoValueTF{#1}%
    {\ZSF@writeidxsee{#2}{#3}}%
    {\ZSF@writeidxsee{#1@#2}{#3}}}

% \ZSFkeyword (Basisdefinition in 50_typography_semantics) indexiert
% jetzt automatisch mit — Rendering bleibt \textbf.
\RenewDocumentCommand{\ZSFkeyword}{s o m}{%
  \textbf{#3}%
  \IfBooleanF{#1}{%
    \IfNoValueTF{#2}{\ZSFindex{#3}}{\ZSFindex{#2}}}}
\makeatother

% ------------------------------------------------------------
% Rendering der Eintragsnummern (Encap)
% ------------------------------------------------------------
% makeindex emittiert \ZSFidxnum{6.2} bzw. \ZSFidxnum{6.2.1}; im
% Seitenzahl-Modus wird sie als letzte Komponente angehängt. Die
% Kapitelziffer-Extraktion funktioniert wie bei
% \ZSFref (\ZSF@refExtractChap aus 50_typography_semantics). Die Farbwahl
% laeuft durch denselben Modulo-Lookup wie \StartChapter.
\makeatletter
\newcommand{\ZSF@idxrendersection}[1]{%
  \edef\ZSF@idxTargetChap{\expandafter\ZSF@refExtractChap#1.\@nil}%
  \ZSFSetPaletteColorNameFromNumber{\ZSF@idxTargetChap}{\ZSF@idxTargetColor}%
  {\upshape\bfseries\color{\ZSF@idxTargetColor}#1}}
\ExplSyntaxOn
\seq_new:N \l__zsf_index_locator_seq
\tl_new:N \l__zsf_index_section_tl
\tl_new:N \l__zsf_index_page_tl
\tl_new:N \l__zsf_index_pending_section_tl
\tl_new:N \l__zsf_index_pending_page_tl
\bool_new:N \l__zsf_index_locator_multiple_bool
\cs_new_protected:Npn \ZSFIndexLocatorListStart
  {\tl_clear:N \l__zsf_index_pending_section_tl \tl_clear:N \l__zsf_index_pending_page_tl \bool_set_false:N \l__zsf_index_locator_multiple_bool}
\cs_new_protected:Npn \ZSFIndexLocatorListMultiple
  {
    \bool_set_true:N \l__zsf_index_locator_multiple_bool
    \tl_if_empty:NF \l__zsf_index_pending_section_tl
      {\exp_args:NV \ZSF@idxrendersection \l__zsf_index_pending_section_tl}
    \tl_clear:N \l__zsf_index_pending_section_tl
    \tl_clear:N \l__zsf_index_pending_page_tl
  }
\cs_new_protected:Npn \ZSFIndexLocatorListFinish
  {
    \bool_if:NF \l__zsf_index_locator_multiple_bool
      {\tl_if_empty:NF \l__zsf_index_pending_section_tl
        {\mbox{\exp_args:NV \ZSF@idxrendersection \l__zsf_index_pending_section_tl\hspace{.2em}\textperiodcentered\hspace{.2em}(S.~\tl_use:N \l__zsf_index_pending_page_tl)}}}
    \tl_clear:N \l__zsf_index_pending_section_tl
    \tl_clear:N \l__zsf_index_pending_page_tl
  }
\cs_new_protected:Npn \ZSFIndexRenderLocator #1
  {
    \seq_set_split:Nnn \l__zsf_index_locator_seq {.} {#1}
    \seq_pop_right:NN \l__zsf_index_locator_seq \l__zsf_index_page_tl
    \tl_set:Nx \l__zsf_index_section_tl
      {\seq_use:Nn \l__zsf_index_locator_seq {.}}
    \bool_if:NTF \l__zsf_index_locator_multiple_bool
      {\exp_args:NV \ZSF@idxrendersection \l__zsf_index_section_tl}
      {\tl_set:Nx \l__zsf_index_pending_section_tl {\tl_use:N \l__zsf_index_section_tl}\tl_set:Nx \l__zsf_index_pending_page_tl {\tl_use:N \l__zsf_index_page_tl}}
  }
\ExplSyntaxOff
\newcommand{\ZSFidxnum}[1]{%
  \ifZSFIndexShowPageNumber
    \ZSFIndexRenderLocator{#1}%
  \else
    \ZSF@idxrendersection{#1}%
  \fi}
\newcommand{\ZSFidxsee}[1]{\ZSF@idxrendersection{#1}}
\makeatother

% see-Verweise zeigen zusätzlich die (farbcodierte) Abschnittsnummer des
% Ziels an, damit ein Synonym-Lookup in EINEM Schritt landet — ohne diese
% Zeile griffe makeidx' Default \see, der die Nummer (#2) verwirft.
% makeindex liefert \see{Ziel}{Nummer}; #2 ist die am kanonischen Ort
% eingefrorene Zielnummer (Konvention: \ZSFindexsee am Zielort platzieren).
\renewcommand*{\see}[2]{\emph{\seename} #1\enspace\ZSFidxsee{#2}}

% ------------------------------------------------------------
% theindex-Umgebung (Override) + Buchstaben-Heading
% ------------------------------------------------------------
% article-Default macht \twocolumn — illegal in multicols*. Stattdessen
% einfache hängende Liste im laufenden 4-Spalten-Fluss.
\makeatletter
\newcommand{\ZSF@idxitem}{\par\hangindent 1.2em\relax}
\newcommand{\ZSF@idxsubitem}{\par\hangindent 2em\hspace*{1em}}
\newcommand{\ZSF@idxsubsubitem}{\par\hangindent 2.6em\hspace*{1.8em}}
\newcommand{\ZSF@idxspace}{\par\addvspace{\ZSFspaceM}}
\renewenvironment{theindex}{%
  \begingroup
  \RaggedRight
  \setlength{\parindent}{0pt}%
  \setlength{\parskip}{0pt plus 0.3pt}%
  \let\item\ZSF@idxitem
  \let\subitem\ZSF@idxsubitem
  \let\subsubitem\ZSF@idxsubsubitem
  \let\indexspace\ZSF@idxspace
}{\par\endgroup}
\makeatother
\newcommand{\ZSFIndexLetter}[1]{%
  \par\addvspace{\ZSFspaceM}%
  {\ZSFfontBoxTitle #1}\par\nobreak\addvspace{\ZSFspaceXS}}
